\begin{tikzpicture}[
  font=\scriptsize,
  execute at begin node={\hyphenpenalty=10000\exhyphenpenalty=10000},
  aktor/.style={draw, rounded corners, fill=blue!8, minimum width=2.1cm, minimum height=0.8cm, align=center},
  garis/.style={dashed, gray},
  ket/.style={font=\tiny, align=center},
  tekan/.style={font=\tiny, align=center, fill=orange!12, inner sep=1.5pt, rounded corners=1pt},
  panah/.style={-{Stealth[length=2mm]}, thick},
  balik/.style={-{Stealth[length=2mm]}, thick, dashed}
]
  \node[aktor] (pasien) at (0,0) {Pasien};
  \node[aktor] (aplikasi) at (3.1,0) {Medikara+\\{\tiny\mdseries klien publik}};
  \node[aktor] (sa) at (6.8,0) {Server otorisasi};
  \node[aktor] (api) at (9.9,0) {API janji temu\\{\tiny\mdseries server sumber daya}};
  \draw[garis] (pasien.south) -- (0,-5.05);
  \draw[garis] (aplikasi.south) -- (3.1,-5.05);
  \draw[garis] (sa.south) -- (6.8,-5.05);
  \draw[garis] (api.south) -- (9.9,-5.05);

  \draw[panah] (0,-0.85) -- node[above, ket] {1. membuka fitur} (3.1,-0.85);
  \node[tekan] at (4.95,-1.2) {2. buat \texttt{code\_verifier},\\hitung \texttt{code\_challenge}};
  \draw[panah] (3.1,-1.8) -- node[above, tekan] {3. pengalihan, \texttt{code\_challenge}} (6.8,-1.8);
  \draw[panah] (0,-2.3) -- node[above, ket, pos=0.75] {4. masuk dan menyetujui} (6.8,-2.3);
  \draw[balik] (6.8,-2.8) -- node[above, ket] {5. kode otorisasi} (3.1,-2.8);
  \draw[panah] (3.1,-3.3) -- node[above, tekan] {6. kode dan \texttt{code\_verifier}} (6.8,-3.3);
  \draw[balik] (6.8,-3.8) -- node[above, ket] {7. hash cocok: token akses} (3.1,-3.8);
  \draw[panah] (3.1,-4.3) -- node[above, ket, pos=0.3] {8. permintaan dengan token} (9.9,-4.3);
  \draw[balik] (9.9,-4.8) -- node[above, ket, pos=0.3] {9. periksa hak atas objek: data atau 403} (3.1,-4.8);

  \node[ket] at (4.95,-5.3) {Pihak yang hanya menangkap kode pada langkah 5 tidak dapat menebusnya, karena tidak memegang \texttt{code\_verifier}};
\end{tikzpicture}
